% cheatsheet/cheatsheet-body.tex
% Two-column A4 cheat-sheet, single source for student + instructor builds.
\documentclass[a4paper,9pt,landscape]{extarticle}

% version-flags.tex
% Single-source instructor/student gating. The Makefile invokes pdflatex with
%   \def\instructorversion{} \input{<artifact>.tex}
% for the instructor build, and a plain include for the student build.
% Each artifact's source begins with % version-flags.tex
% Single-source instructor/student gating. The Makefile invokes pdflatex with
%   \def\instructorversion{} \input{<artifact>.tex}
% for the instructor build, and a plain include for the student build.
% Each artifact's source begins with % version-flags.tex
% Single-source instructor/student gating. The Makefile invokes pdflatex with
%   \def\instructorversion{} \input{<artifact>.tex}
% for the instructor build, and a plain include for the student build.
% Each artifact's source begins with \input{../shared/version-flags.tex}.

\newif\ifinstructor
\newif\ifstudent

% Default: student version. Instructor build sets \instructorversion via -def.
\ifdefined\instructorversion
  \instructortrue
  \studentfalse
\else
  \instructorfalse
  \studenttrue
\fi
.

\newif\ifinstructor
\newif\ifstudent

% Default: student version. Instructor build sets \instructorversion via -def.
\ifdefined\instructorversion
  \instructortrue
  \studentfalse
\else
  \instructorfalse
  \studenttrue
\fi
.

\newif\ifinstructor
\newif\ifstudent

% Default: student version. Instructor build sets \instructorversion via -def.
\ifdefined\instructorversion
  \instructortrue
  \studentfalse
\else
  \instructorfalse
  \studenttrue
\fi


\usepackage[T1]{fontenc}
\usepackage[utf8]{inputenc}
\usepackage{lmodern}
\usepackage[a4paper, landscape, margin=1.2cm]{geometry}
\usepackage{multicol}
\usepackage{enumitem}
\setlist[itemize]{itemsep=1pt, topsep=2pt, leftmargin=*, parsep=0pt}
\setlist[enumerate]{itemsep=1pt, topsep=2pt, leftmargin=*, parsep=0pt}
\usepackage{xcolor}
\usepackage{hyperref}
\hypersetup{colorlinks=true, urlcolor=blue!60!black, pdftitle={Introduction to ArduCopter --- Cheat-sheet}}
\usepackage{url}
\usepackage{titlesec}
\titleformat{\section}{\large\bfseries}{}{0pt}{}
\titleformat{\subsection}{\bfseries}{}{0pt}{}
\titlespacing*{\section}{0pt}{4pt}{2pt}
\titlespacing*{\subsection}{0pt}{3pt}{1pt}
\setlength{\columnsep}{1cm}
\setlength{\parindent}{0pt}
\pagestyle{empty}

% macros.tex
% Citation hyperlinks point at this fork on the GNC-0.1 branch (where the course tree lives).
% Display text follows citation-rigor.md: "path:line" or "path:start-end".
% Link target is the GitHub blob URL on the project's actual remote.

\providecommand{\repourl}{https://github.com/hgrobotics/ardupilot_course/blob/GNC-0.1}

% \citeline{path}{line} — single-line cite (e.g. \citeline{ArduCopter/Copter.h}{181})
\newcommand{\citeline}[2]{\href{\repourl/#1\#L#2}{\nolinkurl{#1}:#2}}

% \citerange{path}{start}{end} — line-range cite
\newcommand{\citerange}[3]{\href{\repourl/#1\#L#2-L#3}{\nolinkurl{#1}:#2-#3}}

% \citefile{path} — whole-file cite (no line)
\newcommand{\citefile}[1]{\href{\repourl/#1}{\nolinkurl{#1}}}

% Sibling-course pointers (relative paths inside this repo, used in instructor "advanced pointers" notes).
\newcommand{\siblingplane}{\href{\repourl/course/custom_gnc_course_plane.md}{course/custom\_gnc\_course\_plane.md}}
\newcommand{\siblingquadplane}{\href{\repourl/course/custom_gnc_course_quadplane.md}{course/custom\_gnc\_course\_quadplane.md}}

% Copyright line — inserted on title pages / cheatsheet header.
\newcommand{\copyrightline}{\copyright\ 2026 HG Robotics Co., Ltd.}

% Inline param-name and code-path styling (no inline code listings on slides per pedagogy).
\newcommand{\param}[1]{\texttt{#1}}
\newcommand{\codepath}[1]{\texttt{\nolinkurl{#1}}}
\newcommand{\mavcmd}[1]{\texttt{#1}}
% Note: \mode is already defined by Beamer (for overlay specs). Use \flmode (flight mode).
\newcommand{\flmode}[1]{\textsc{#1}}

% Instructor-note environments. Suppressed in student build via the \ifinstructor gate.
% Used for: (1) expected answers + timing, (2) anticipated student FAQs, (3) speaker notes,
% (4) advanced-material pointers.
\newcommand{\instnote}[1]{\ifinstructor\begin{block}{Instructor note}#1\end{block}\fi}
\newcommand{\insttiming}[1]{\ifinstructor\begin{block}{Pacing}#1\end{block}\fi}
\newcommand{\instfaq}[2]{\ifinstructor\begin{block}{Anticipated Q: #1}#2\end{block}\fi}
\newcommand{\instadvanced}[1]{\ifinstructor\begin{block}{Advanced material pointer}#1\end{block}\fi}


% Cheat-sheet doesn't carry full instructor blocks; instructor edition adds answer-key
% callouts inline only.
\renewcommand{\instnote}[1]{\ifinstructor\par\smallskip\textit{\footnotesize Instructor: #1}\par\fi}
\renewcommand{\insttiming}[1]{}
\renewcommand{\instfaq}[2]{}
\renewcommand{\instadvanced}[1]{\ifinstructor\par\smallskip\textit{\footnotesize Advanced: #1}\par\fi}

\begin{document}
\begin{center}
  \textbf{\large Introduction to ArduCopter --- Cheat-sheet}\quad\textbar\quad
  \ifinstructor Instructor edition\else Student edition\fi\quad\textbar\quad
  Branch \texttt{GNC-0.1}\quad\textbar\quad
  \footnotesize\copyrightline
\end{center}
\hrule\medskip

\begin{multicols*}{2}

\section*{MAVProxy commands used in this course}
Only these. No others.
\begin{itemize}
  \item \texttt{mode <NAME>} --- request a flight-mode change.
  \item \texttt{arm throttle} --- arm motors.
  \item \texttt{rc 3 1700} --- throttle stick to 1700~$\mu$s PWM.
  \item \texttt{takeoff <m>} --- guided take-off to altitude (Step A).
  \item \texttt{param show <NAME>} --- read a parameter.
  \item \texttt{param set <NAME> <VAL>} --- write a parameter.
\end{itemize}

\section*{SITL launch}
\begin{itemize}
  \item \texttt{./waf configure --board sitl}
  \item \texttt{./waf copter}
  \item \texttt{Tools/autotest/sim\_vehicle.py -v ArduCopter --console --map}
\end{itemize}
Three windows: console (status), command (you type), map.\\
\textbf{Pass:} \texttt{HEARTBEAT} + map render + \texttt{online system 1} within $\sim$30~s.

\section*{Flight modes --- 3 buckets}
\begin{itemize}
  \item \textbf{Manual} --- \flmode{STABILIZE}, \flmode{ACRO}.
  \item \textbf{Stabilized-with-altitude} --- \flmode{ALT\_HOLD}.
  \item \textbf{Autonomous} --- \flmode{AUTO}, \flmode{GUIDED}, \flmode{RTL}, \flmode{LAND}.
\end{itemize}
Mode IDs (subset): \flmode{STABILIZE}=0, \flmode{ALT\_HOLD}=2, \flmode{LAND}=9.\\
Mode entries: \citerange{ArduCopter/mode.h}{77}{109}.

\section*{Mode-switch refusal}
\begin{itemize}
  \item \texttt{requires position} $\rightarrow$ \citeline{ArduCopter/mode.cpp}{394}.
  \item \texttt{need alt estimate} $\rightarrow$ \citeline{ArduCopter/mode.cpp}{404}.
\end{itemize}
Function: \texttt{Copter::set\_mode}, \citerange{ArduCopter/mode.cpp}{313}{396}.

\section*{The init / run pattern}
Every mode has \texttt{init()} (once on entry) and \texttt{run()} (high rate).
\begin{itemize}
  \item \citerange{ArduCopter/mode_althold.cpp}{9}{22} --- \texttt{ModeAltHold::init}.
  \item \citerange{ArduCopter/mode_althold.cpp}{26}{104} --- \texttt{ModeAltHold::run}.
  \item \citerange{ArduCopter/mode_stabilize.cpp}{9}{64} --- \texttt{ModeStabilize::run}.
\end{itemize}

\section*{First flight (Lab L2)}
\begin{enumerate}
  \item \texttt{mode STABILIZE}
  \item \texttt{arm throttle}
  \item \texttt{rc 3 1700}
  \item Wait altitude $> 10$~m.
  \item \texttt{mode LAND}
  \item Wait \texttt{Disarming motors}.
\end{enumerate}
\textbf{Pass:} \texttt{Disarming motors} within 90~s of arm.

\section*{Parameters}
Persistent runtime config; not compiled in. Annotation-block example at \citerange{ArduCopter/Parameters.cpp}{44}{51} (\texttt{PILOT\_THR\_FILT}); keys: \texttt{@Param}, \texttt{@DisplayName}, \texttt{@Description}, \texttt{@Range}/\texttt{@Values}, \texttt{@Units}, \texttt{@User}.\\
Real table opening: \citerange{ArduCopter/Parameters.cpp}{33}{67}.\\
Mode bands: \citerange{ArduCopter/Parameters.cpp}{149}{191} (\param{FLTMODE1}--\param{FLTMODE6}, \param{FLTMODE\_CH}).

\columnbreak

\section*{Sensors $\rightarrow$ EKF $\rightarrow$ scheduler}
\begin{itemize}
  \item Sensors: IMU, baro, mag, GPS.
  \item EKF: estimator. Black box for this course.
  \item Scheduler: \texttt{FAST\_TASK} (every loop) or \texttt{SCHED\_TASK(func, rate, max\_us, prio)}.
\end{itemize}
Anchors:
\begin{itemize}
  \item \citerange{ArduCopter/Copter.cpp}{113}{149} --- \texttt{scheduler\_tasks[]} opening.
  \item \citeline{ArduCopter/Copter.cpp}{117} --- \texttt{run\_rate\_controller\_main}.
  \item \citeline{ArduCopter/Copter.cpp}{127} --- \texttt{read\_AHRS} (= EKF).
  \item \citeline{ArduCopter/Copter.cpp}{201} --- \texttt{ekf\_check, 10~Hz}.
  \item \citeline{ArduCopter/Copter.cpp}{998} --- entry point macro.
  \item \citerange{libraries/AP_Vehicle/AP_Vehicle.cpp}{558}{566} --- main loop $=$ \texttt{scheduler.loop()}.
\end{itemize}
\texttt{SCHED\_LOOP\_RATE} default: 400~Hz.

\section*{The data-path sentence (memorise)}
\begin{quote}\small
Stick $\rightarrow$ mode \texttt{run()} $\rightarrow$ attitude controller $\rightarrow$ rate controller $\rightarrow$ PID $\rightarrow$ motor mixer $\rightarrow$ ESC.
\end{quote}
Stage 3 anchors:
\begin{itemize}
  \item \citerange{ArduCopter/Attitude.cpp}{10}{24} --- \texttt{run\_rate\_controller\_main}.
  \item \citerange{libraries/AC_AttitudeControl/AC_AttitudeControl_Multi.cpp}{457}{485} --- body.
  \item Roll/pitch/yaw PIDs: \citeline{libraries/AC_AttitudeControl/AC_AttitudeControl_Multi.cpp}{473}, \citeline{libraries/AC_AttitudeControl/AC_AttitudeControl_Multi.cpp}{476}, \citeline{libraries/AC_AttitudeControl/AC_AttitudeControl_Multi.cpp}{479}.
\end{itemize}
PID: \citerange{libraries/AC_PID/AC_PID.cpp}{196}{272}. Identify \texttt{P\_out = (\_error * \_kp);} and stop.\\
Mixer opening: \citerange{libraries/AP_Motors/AP_MotorsMatrix.cpp}{213}{244}.\\
X-frame angles: \citerange{libraries/AP_Motors/AP_MotorsMatrix.cpp}{592}{602} ($+45^\circ, -135^\circ, -45^\circ, +135^\circ$, alternating CCW/CW).

\section*{Failsafes}
Four to recognise: RC loss, battery, GCS link loss, \textbf{EKF variance}.\\
EKF-failsafe code path: \citerange{ArduCopter/ekf_check.cpp}{30}{90}.\\
Names: \texttt{ekf\_check\_state.fail\_count}, \texttt{.bad\_variance}, \texttt{g.fs\_ekf\_thresh} ($=$ \param{FS\_EKF\_THRESH}), \texttt{EKF\_CHECK\_ITERATIONS\_MAX}.\\
Three trigger lines:
\begin{itemize}
  \item \citeline{ArduCopter/ekf_check.cpp}{83} --- \texttt{LOGGER\_WRITE\_ERROR(EKFCHECK, BAD\_VARIANCE)}.
  \item \citeline{ArduCopter/ekf_check.cpp}{86} --- \texttt{EKF variance:} \texttt{CRITICAL} statustext.
  \item \citeline{ArduCopter/ekf_check.cpp}{89} --- \texttt{failsafe\_ekf\_event()}.
\end{itemize}
Pre-arm posture (same shape, at boot): \citerange{ArduCopter/AP_Arming_Copter.cpp}{8}{20}.

\section*{Closing lab (L3) --- Step A clean run}
\begin{verbatim}
arm throttle
mode GUIDED
takeoff 30
(hover 30 s)
mode RTL
(wait disarm)
\end{verbatim}
\textbf{Pass:} disarm within 180~s, no \texttt{CRITICAL} statustext.

\section*{Closing lab (L3) --- Step B failsafe}
At t = 60~s after arm: \texttt{param set SIM\_GPS\_DISABLE 1}.\\
\textbf{Pass} (all 4):
\begin{enumerate}
  \item \texttt{EKF variance:} \texttt{CRITICAL} within 30~s of injection.
  \item Mode change to \flmode{LAND}/\flmode{RTL} within 30~s.
  \item Disarm within 240~s of original arm.
  \item Dataflash \texttt{ERR}: \texttt{Subsys=EKFCHECK}, \texttt{ECode=BAD\_VARIANCE} via \texttt{mavlogdump.py --types ERR}.
\end{enumerate}
Fingerprint: \citerange{ArduCopter/ekf_check.cpp}{79}{89}.

\ifinstructor
\section*{Instructor extras}
\begin{itemize}
  \item Pacing: Day 1 hands-on $\sim$59\%; Day 2 hands-on $\sim$41\% (load-bearing on Module 2.4).
  \item Don't unpack PID, EKF, or spool-state internals --- those are the GNC course's job.
  \item Advanced pointers for Q\&A: \siblingplane, \siblingquadplane.
  \item Apt-mirror diagnostic: verify \texttt{apt-cache madison <pkg>} against \texttt{archive.ubuntu.com}; switch the configured mirror if it is behind.
\end{itemize}
\fi

\end{multicols*}
\end{document}
